
\documentclass[11pt]{article}

\usepackage{amsmath,amsthm,amssymb,mathtools,fullpage,enumitem}

\newtheorem{theorem}{Theorem}
\newtheorem{lemma}{Lemma}
\newtheorem{proposition}{Proposition}
\newtheorem{remark}{Remark}

\newcommand{\R}{\mathbb R}
\newcommand{\eps}{\varepsilon}
\newcommand{\Om}{\Omega}
\newcommand{\deliso}{\delta_{\mathrm{iso}}}

\title{Notes on Quantitative Stability for Almost Serrin Domains}
\author{}
\date{}

\begin{document}

\maketitle

\section{Setup}

Let $\Omega \subset \mathbb R^n$ be a bounded smooth domain and let
$u$ solve
\[
-\Delta u = 1 \quad \text{in }\Omega,
\qquad
u=0 \quad \text{on }\partial\Omega.
\]

Assume that $\Omega$ is almost isoperimetric, namely
\[
\deliso(\Omega)
:=
\frac{P(\Omega)}
{n\omega_n^{1/n}|\Omega|^{(n-1)/n}}
-1
\ll 1,
\]
and assume that the Serrin deficit is small:
\[
\eps^2
:=
\int_{\partial\Omega}
\left(
|\nabla u|-c_\Omega
\right)^2
\,d\mathcal H^{n-1},
\]
where
\[
c_\Omega
=
\frac{|\Omega|}{P(\Omega)}
=
\fint_{\partial\Omega} |\nabla u|.
\]

The goal is to show that $\partial\Omega$ is close to a sphere,
eventually as a radial graph.

\section{Known quantitative ingredients}

\subsection{Quantitative isoperimetry}

By the sharp quantitative isoperimetric inequality of
Fusco--Maggi--Pratelli,
there exists a ball $B_R(x_0)$ such that
\[
|\Omega \Delta B_R(x_0)|
\le
C(n)|\Omega|\deliso(\Omega)^{1/2}.
\]

Thus small isoperimetric deficit implies closeness in measure to a ball.

\subsection{Quantitative Serrin stability}

Classical quantitative Serrin stability results
(Magnanini--Poggesi and related works)
show that if
\[
\left\|
|\nabla u|-c_\Omega
\right\|_{L^2(\partial\Omega)}
\]
is small, then $\Omega$ is trapped inside a thin annulus:
\[
B_{\rho_i}(z)\subset \Omega \subset B_{\rho_e}(z),
\]
with
\[
\rho_e-\rho_i
\le
C \eps^\theta,
\]
provided one assumes geometric nondegeneracy
(interior/exterior sphere condition, $C^{2,\alpha}$ bounds, etc.).

The discussion below aims at replacing those geometric constants
by the pair
\[
(\deliso,\eps).
\]

\section{Flux identities}

Fix $x\in \partial\Omega$ and define
\[
A_r:=\Omega\cap B_r(x),
\]
\[
\Gamma_r:=\partial\Omega\cap B_r(x),
\]
\[
\Sigma_r:=\Omega\cap \partial B_r(x).
\]

Integrating $-\Delta u=1$ over $A_r$ gives
\[
|A_r|
=
\int_{\Gamma_r} |\nabla u|
-
\int_{\Sigma_r} \partial_{\nu_{B_r}}u.
\tag{1}
\]

Define the localized deficit
\[
\eps_r^2
:=
\int_{\Gamma_r}
(|\nabla u|-c_\Omega)^2.
\]

By Cauchy--Schwarz,
\[
\int_{\Gamma_r} |\nabla u|
\ge
c_\Omega \mathcal H^{n-1}(\Gamma_r)
-
\eps_r
\mathcal H^{n-1}(\Gamma_r)^{1/2}.
\]

Hence from (1),
\[
\eps_r
\mathcal H^{n-1}(\Gamma_r)^{1/2}
\ge
c_\Omega \mathcal H^{n-1}(\Gamma_r)
-|A_r|
-
\left|
\int_{\Sigma_r}\partial_{\nu_{B_r}}u
\right|.
\tag{2}
\]

\section{Control of the spherical flux}

Using comparison with balls, one has the global gradient estimate
\[
\|\nabla u\|_{L^\infty(\Omega)}
\le
C(n,\operatorname{diam}\Omega).
\]

By coarea,
\[
\int_r^{2r}
\mathcal H^{n-1}(\Sigma_s)\,ds
=
|\Omega\cap(B_{2r}(x)\setminus B_r(x))|.
\]

Thus if
\[
|\Omega\cap B_{2r}(x)|
\le
\eta r^n,
\]
then there exists $s\in(r,2r)$ such that
\[
\mathcal H^{n-1}(\Sigma_s)
\le
C\eta r^{n-1}.
\]

Consequently,
\[
\left|
\int_{\Sigma_s}\partial_{\nu_{B_s}}u
\right|
\le
C(n,\operatorname{diam}\Omega)\eta r^{n-1}.
\tag{3}
\]

\section{Relative isoperimetric inequality}

The relative isoperimetric inequality in balls states that
\[
\min\{
|\Omega\cap B_r(x)|,
|B_r(x)\setminus \Omega|
\}^{\frac{n-1}{n}}
\le
C_n P(\Omega;B_r(x)).
\]

Therefore small volume fraction implies a lower bound
for the local perimeter:
\[
P(\Omega;B_r(x))
\ge
c_n |\Omega\cap B_r(x)|^{(n-1)/n}.
\]

This removes the need to assume a priori lower perimeter density.

\section{Exclusion of macroscopic tentacles}

Suppose that
\[
|\Omega\cap B_{2r}(x)|
\le
\eta r^n.
\]

Combining the local perimeter lower bound,
the flux estimate (3),
and inequality (2),
one obtains
\[
\eps_r^2
\gtrsim
c_\Omega^2
P(\Omega;B_r(x)).
\]

Using relative isoperimetry,
\[
\eps_r^2
\gtrsim
|\Omega\cap B_r(x)|^{(n-1)/n}.
\]

Thus bad interior density forces positive Serrin deficit.

Equivalently,
if the global Serrin deficit satisfies
\[
\eps^2
=
\int_{\partial\Omega}
(|\nabla u|-c_\Omega)^2
\ll1,
\]
then bad density can occur only below the scale
\[
r_*
\sim
\eps^{2/(n-1)}.
\]

More precisely, for every $x\in\partial\Omega$ and every
\[
r\gg \eps^{2/(n-1)},
\]
one has the lower density estimate
\[
|\Omega\cap B_r(x)|
\ge
c r^n.
\]

\section{Exterior density}

Quantitative isoperimetry gives
\[
|\Omega\Delta B_R(x_0)|
\lesssim
\deliso(\Omega)^{1/2}.
\]

Thus large exterior holes or fjords would either

\begin{enumerate}[label=(\roman*)]
\item contribute substantial volume discrepancy against the nearby ball,
or
\item produce a thin geometric region whose flux behavior forces
nontrivial Serrin deficit.
\end{enumerate}

Hence one expects the two-sided density estimate
\[
c r^n
\le
|\Omega\cap B_r(x)|
\le
(1-c)r^n
\]
for all boundary points $x$ and all scales
\[
r
\ge
C
\left(
\eps^{2/(n-1)}
+
\deliso(\Omega)^{1/(2n)}
\right).
\]

\section{Adapted Magnanini--Poggesi step}

Define
\[
q(x)
=
\frac{R^2-|x-z|^2}{2n},
\]
where $z$ is a suitable center point, and let
\[
h:=q-u.
\]

Since
\[
-\Delta q=1,
\]
the function $h$ is harmonic in $\Omega$.

Moreover,
\[
h|_{\partial\Omega}
=
\frac{R^2-|x-z|^2}{2n}.
\]

Hence if
\[
\rho_i
=
\min_{x\in\partial\Omega}|x-z|,
\qquad
\rho_e
=
\max_{x\in\partial\Omega}|x-z|,
\]
then
\[
\rho_e^2-\rho_i^2
=
2n\,
\operatorname{osc}_{\partial\Omega} h.
\]

The classical Magnanini--Poggesi argument estimates
$\operatorname{osc}_{\partial\Omega} h$
through weighted Hessian control:
\[
\int_\Omega
u |\nabla^2 h|^2
\le
C
\int_{\partial\Omega}
(|\nabla u|-c_\Omega)^2.
\]

In the present setting, the proposed replacement estimate is
\[
\operatorname{osc}_{\partial\Omega} h
\le
C
\|\nabla^2 h\|_{L^2(\Omega,u\,dx)}^\theta
+
C r_*,
\]
where
\[
r_*
=
C
\left(
\eps^{2/(n-1)}
+
\deliso(\Omega)^{1/(2n)}
\right).
\]

Combining these estimates yields
\[
\rho_e-\rho_i
\le
C
\left(
\eps^\theta
+
r_*
\right).
\]

Thus the domain is trapped in a thin annulus with constants depending only on
\[
n,
\quad
|\Omega|,
\quad
\operatorname{diam}\Omega,
\]
and the small deficits.

\section{Remaining steps}

The discussion above leaves the following technical points open:

\begin{enumerate}[label=(\arabic*)]
\item Establishing the adapted harmonic interpolation estimate rigorously under two-sided density assumptions above scale $r_*$.

\item Showing that thin annular trapping plus small isoperimetric deficit imply that $\partial\Omega$ is a single radial graph.

\item Upgrading the radial graph to $C^{1,\alpha}$ regularity.
\end{enumerate}

These appear plausible within the framework developed above,
but they still require a complete proof.

\end{document}
